\subsection{DECOMPOSITION OF A HOMOMORPHISM}

PROPOSITION 6. Let G be a group with operators and G' a magma with an action by Ω, written exponentially. Let f: G → G' be a homomorphism of the magna G into the magma G' such that, for all α \ensuremath{\in} Ω and all x \ensuremath{\in} G, \(f(x^{\alpha}) = f(x)^{\alpha}\) . Then \(f(\mathbf{G})\) is a stable subset of G' under the law on G' and the action of Ω; the set \(f(\mathbf{G})\) with the induced laws is a group with operators and the mapping \(x \mapsto f(x)\) of G into \(f(\mathbf{G})\) is a homomorphism of groups with operators.

By virtue of § 1, no. 4, Proposition 1, \(f(\mathbf{G})\) is a stable subset of \(\mathbf{G}'\) under the internal law on \(\mathbf{G}'\) . For every element \(x \in \mathbf{G}\) and for every operator \(\alpha\) , \(f(x)^{\alpha} = f(x^{\alpha}) \in f(\mathbf{G})\) and therefore \(f(\mathbf{G})\) is stable under the action of \(\Omega\) on \(\mathbf{G}'\) . Writing the internal law of \(\mathbf{G}'\) multiplicatively,

\[
\begin{array}{r l} (f (x) f (y)) f (z) & = f (x y) f (z) = f ((x y) z) = f (x (y z)) = f (x) f (y z) \\ & \qquad \qquad \qquad \qquad \qquad \qquad \qquad \qquad \qquad \qquad = f (x) (f (y) f (z)) \end{array}
\]

for all elements \(x, y, z\) in \(\mathbf{G}\) ; therefore the induced law on \(f(\mathbf{G})\) is associative.

Let \(e\) be the identity element of \(G\) . Its image \(f(e)\) is an identity element of \(f(G)\) (§ 2, no. 1). Every element of \(f(G)\) is invertible in \(f(G)\) (§ 2, no. 3). Therefore the law induced on \(f(G)\) by the internal law on \(G'\) is a group law. For all elements \(x\) and \(y\) in \(G\) and every operator \(\alpha\) ,

\[
(f (x) f (y)) ^ {\alpha} = (f (x y)) ^ {\alpha} = f ((x y) ^ {\alpha}) = f \left(x ^ {\alpha} y ^ {\alpha}\right) = f \left(x ^ {\alpha}\right) f \left(y ^ {\alpha}\right) = (f (x)) ^ {\alpha} (f (y)) ^ {\alpha}
\]

which shows that the action of \(\Omega\) is distributive with respect to the group law on \(f(\mathrm{G})\) . Therefore \(f(\mathrm{G})\) with the induced laws is a group with operators and clearly the mapping \(x \mapsto f(x)\) is a homomorphism of groups with operators.

DEFINITION 8. Let \(f: \mathbf{G} \to \mathbf{G}'\) be a homomorphism of groups with operators. The inverse image of the identity element of \(\mathbf{G}'\) is called the kernel of \(f\) .

The kernel of \(f\) is often denoted by \(\operatorname{Ker}(f)\) and the image \(f(\mathbf{G})\) of \(f\) is sometimes denoted by \(\operatorname{Im}(f)\) .

THEOREM 3. Let \(f: \mathbf{G} \to \mathbf{G}'\) be a homomorphism of groups with operators.

\begin{enumerate}
\def\labelenumi{(\alph{enumi})}
\item
  \(\operatorname{Ker}(f)\) is a normal stable subgroup of \(\mathbf{G}\) ;
\item
  \(\operatorname{Im}(f)\) is a stable subgroup of \(\mathbf{G}'\) ;
\item
  the mapping \(f\) is compatible with the equivalence relation defined on \(\mathbf{G}\) by \(\operatorname{Ker}(f)\)
\item
  the mapping \(\tilde{f}\colon \mathbf{G} / \mathbf{Ker}(f)\to \operatorname {Im}(f)\) derived from \(f\) by passing to the quotient is an isomorphism of groups with operators;
\item
  \(f = \iota \circ \tilde{f} \circ \pi\) , where \(\iota\) is the canonical injection of \(\operatorname{Im}(f)\) into \(G'\) and \(\pi\) is the canonical homomorphism of \(G\) onto \(G / \operatorname{Ker}(f)\) .
\end{enumerate}

Assertion (b) follows from Proposition 6. The equivalence relation \(f(x) = f(y)\) on \(G\) is compatible with the group with operators structure on \(G\) . By Theorem 2 (no. 4), it is therefore of the form \(y \in xH\) , where \(H\) is a normal stable subgroup of \(G\) and \(H\) is the class of the identity element, whence \(H = \operatorname{Ker}(f)\) . Assertions (a), (c) and (d) then follow. Assertion (e) is obvious (Set Theory, II, § 6, no. 5).

